%% ma: L12-B08-0023
%% lop: 12
%% chuong: 2
%% bai: 8
%% dang: 12.8.15
%% loai: TN4
%% muc: TH
%% dap_an: "A"
%% nguon: "(NBV)-ĐỀ SỐ 16-MỖI NGÀY 1 ĐỀ THI 2025.pdf (D16, MỖI NGÀY 1 ĐỀ - Đề số 16), Phần I Câu 11"
%% kien_thuc: "Bài 8 (biểu thức tọa độ của các phép toán vectơ)"
%% trang_thai: cho-duyet
%% ly_do: "Dạng mới đề xuất 12.8.15 chờ giáo viên duyệt"
%% ghi_chu: "Đáp án gốc A đúng (x = 2, y = 1)."
\begin{ex}
Trong không gian $Oxyz$, cho các vectơ $\vec a=(2;-1;2)$, $\vec b=(-3;1;-1)$ và $\vec c=(1;-1;3)$. Biết rằng $\vec c=x\vec a+y\vec b$ với $x,y$ là các số thực. Khi đó $x+y$ bằng
\choice{\True $3$}{$1$}{$2$}{$4$}
\loigiai{$x\vec a+y\vec b=(2x-3y;\,-x+y;\,2x-y)$. Từ $\vec c=x\vec a+y\vec b$ ta có hệ $\begin{cases}2x-3y=1\\-x+y=-1\\2x-y=3\end{cases}$, suy ra $x=2$, $y=1$ (thỏa cả ba phương trình). Vậy $x+y=3$. Chọn A.}
\end{ex}
